\documentclass{article}
\usepackage[utf8]{inputenc}
\usepackage{hyperref}
\usepackage{fullpage}
\usepackage{amsmath}
\title{subgraph freeness}
\author{talyaa01}
\date{February 2022}

\newcommand{\eps}{\epsilon}
\newcommand{\light}{\textsf{light}}
\newcommand{\heavy}{\textsf{heavy}}
\renewcommand{\ll}{\textsf{$\ell-\ell$}}
\newcommand{\hh}{\textsf{$h-h$}}
\newcommand{\lh}{\textsf{$$\ell-h$}}
\newcommand{\davg}{d_{avg}}
\newcommand{\poly}{\textrm{poly}}

\begin{document}

\iffalse

\maketitle

our jamboard:
\href{https://jamboard.google.com/d/10kNo1nc42vDcbuwcF_MHJMbRpsfyljGIQr6ow0UyukA/viewer?f=0}{jamboard-link}
second board:
\href{https://jamboard.google.com/d/140dJaqwfDGQGzm8S4-iJL_9TREsbfLXTOCvWSCTlU1U/edit?usp=sharing}{second board}

\section{Meeting 8/2/2022}

\subsection{Upper bound ideas for 4-cycle freeness}

\begin{dfn}[light and heavy vertices and edges]
If a vertex $v$ is such that $d(v)\leq \alpha/\eps$ then we say it is \light, and otherwise we say it is \heavy.  We classify edges into three types depending on their endpoints. An edge is denoted $\ll$ if it has two light endpoints, and similarly for $\hh$ and $\lh$,
\end{dfn}


First, we can remove all edges between two heavy vertices (where heavy is a vertex $v$ with  $d(v)>\alpha/\eps$).
We can also ignore vertices that have less than $\eps d(v)$ 4-cycles incident to them.
Procedure:
\begin{enumerate}
    \item \color{blue}{Do we sample vertices uar or do we start from edges and then take their endpoints}

    \item For each sampled vertex $u$:
    \begin{enumerate}
        \item Sample a set $X_u$ of size $O(\sqrt{d(u}/\eps)$ of $u'$s neighbors.
        \item For every sampled light neighbor $v$ of $u$, query \emph{all} neighbors of $v$ - $N_v$
        \item For every two pairs $v,v' \in X_u$, let $\chi_{u,v,v'}=|\Gamma(v)\cap \Gamma(v')|$
    \end{enumerate}

\end{enumerate}
Actually I forgot why we can start with a uniform edge sample and still say that $v$ has $O(\eps d(v))$ neighboring cycles. Don't we lose that when we start with an edge? I think I asked this before, but forgot the answer.

I (Dana) started to answer this question, but then I had a point of difficulty, which I think we may be able to overcome by dealing a bit differently with the issue of vertices that have relatively few (to their degree) 4-cycles residing on them. I'll repeat some things that Talya wrote, and do things slowly and in a bit more detail.

We started with $\Omega(\eps m)$ edge-disjoint $4$-cycles. After marking/removing all H-H edges, we still have $\Omega(\eps m)$ edge-disjoint $4$-cycles with no H-H edge in them. (The reason we might prefer ``mark'' is just so as to remember, that when we sample edges, we might hit such ``unuseful'' H-H edges.)

Next we perform an iterative process in which we allow each vertex $v$ to mark/remove at most $\eps' d(v)$ incident edges that reside on 4-cycles (for $\eps' = \eps/c$ where $c$ is a sufficiently large constant). The point will be that for some vertices we might mark/remove almost all their incident edges, but each vertex will ``pay'' for at most $\eps'$ of this incident edges.

As long as there is a vertex $v$ such that the number of edges incident to $v$ that participate in edge-disjoint cycles (that are still ``alive'', that is with no marked/removed edge) is at most $\eps' d(v)$, we mark/remove all these edges. So when we are done (after removing a total of at most $2\eps' m$ edges) we know we have $\Omega(\eps m)$ edges, that are not H-H and such that each of their endpoints $v$ participates in $\Omega(\eps d(v))$ 4-cycles.

For $\lh$ cycles, the idea was to sample $\sqrt n $ vertices from the $\heavy$ side and $O(\alpha)$ from the light, and then check for all collisions. Here too I forgot why this should work.

\subsection{Meeting 31/3: $O(\sqrt{|H^*|}))$ for $L_H^*-L-H^*$ cycles }
First, we think that the right parameter to work with is $\alpha^*$ (defined in the testing arboricity paper).
I'm actually not sure how this goes technically into the result, but we can think about this later. Perhaps it's similar to "graying out" all edges incident to low degree (less than $\eps \davg$) and low C4 vertices.

Let $\theta$ be some degree threshold, and let $H^*=\{v\mid d(v)\geq \theta\}$. Further let $L=\{v \mid d(v)\leq \alpha/\eps\}$, and $H=V\setminus L\setminus H^*$.
We consider a bipartite graph, where all vertices on the left belong to $L$ and all vertices on the right belong to $H^*$.
Since every vertex in $L$ has degree at least $\eps d(v)$, it holds that each vertex in $H^*$ has at least $\eps \theta \davg$ edge disjoint paths of length to originating in it. Fix  such a vertex $v'$. That is, we have that
\begin{align}\label{eq:sum_l2v'}
\ell_2(v')=\sum_{v\in H^*} \ell_2(v',v) \geq \eps \theta \davg.
\end{align}

Now consider a process that chooses two (edge-disjoint) paths of length two originating in $v'$.
The following holds.
\begin{align} \label{eq:col_prob_l2v}
\Pr[collision] =\sum_{v'\in H^*}\frac{l_2(v,v')}{\ell_2(v)}\cdot \frac{\ell_2(v,v')-1}{\ell_2(v)}=
\sum_{v'\in H^* s.t. {\ell_2(v,v')\geq 2}}\frac{1}{2}\frac{\ell_2(v,v')}{\ell_2(v)}-\frac{|H^*|}{\ell_2(v)}
\geq
\frac{1}{2}\sum_{v'\in H^*}\frac{(\ell_2(v,v'))^2}{(\ell_2(v))^2}-\frac{|H^*|}{(\ell_2(v))^2}
\end{align}
We now lower bound the first term in the RHS, and upper bound the second.
For the first term, by  Cauchy-Schwartz inequality it holds that
\begin{align}\label{eq:lb_first_term}
    \sum_{v\in H^*}\frac{(\ell_2(v',v))^2}{(\ell_2(v'))^2} \geq \sum_{v\in H^*}\frac{|H^*|\left(\frac{(\ell_2(v'))}{|H^*|}\right)^2}{(\ell_2(v'))^2}=\frac{1}{|H^*|}\geq
\end{align}
We turn to upper bound the second term.
By Equation~\ref{eq:sum_l2v'} it holds that $\ell_2(v')\geq \eps \theta \davg$. It also holds that $|H^*|\leq m/\theta$.
Setting $\theta=4\sqrt n/\eps$ ensures the following:
\begin{align}
    |H^*|\leq 2m/\theta,
\end{align}
and
\begin{align}
    \ell_2(v')\geq \eps \davg\theta
\end{align}
Therefore, setting $\theta^2 =4m/\eps\davg \Rightleftarrow \theta=2\sqrt{n/\eps}$, implies that
\begin{align}
    \ell_2(v')\geq 2|H^*|,
\end{align}
implying that
\begin{align}\label{eq:ub_second_term}
    \frac{|H^*|}{(\ell_2(v'))^2)} \leq \frac{1}{4|H^*|}.
\end{align}
Plugging in Equations~\eqref{eq:lb_first_term} and~\eqref{eq:ub_second_term} into the RHS of Equation~\ref{eq:col_prob_l2v}, we get that
\begin{align}
    \sum_{v\in H^*}\frac{(\ell_2(v',v))^2}{(\ell_2(v'))^2} \geq \frac{1}{2|H^*|}.
\end{align}

\subsection{meeting (talya, reut), 13.4 - trying to correct and generalize the $O(\sqrt{|H^*|})$ argument}

As usual, we are only considering the bipartite graph $H^*, L$ and vertices $v$ such that $V$ has at  least $\eps d(v)$ edge participating in $C_4$'s. Let's refer to these edges as red edges. Also, we can only consider edge-disjoint $C_4$'s.

Fix $v\in H^*$. Bucket the vertices of $\Gamma(v)\cap L$ according to their degree, so that $B_i=\{u\in \Gamma(v)\cap L \mid d(u)\in [\eps\davg\cdot (1+\beta)^{i-1}, \eps\davg\cdot (1+\beta)^i) \}$, and there are $O(\log \frac{\alpha}{\eps^2 \davg})=O(\log n)$ buckets.
So there is at least one bucket $i$ so that  $\eps^2 d(v)/\log n$ of $v$'s red edges go to $B^v_i$. Let's denote the (approximate) degree of the vertices in $B^v_i$ by $y$.
We have
\begin{align}
    \Pr[collision] \geq \Pr[\text{collision through }$B_i$]= \sum_{v'\in H^*} \frac{(\ell^{B_i}_2(v,v'))(\ell_2^{B_i}(v,v')-1)}{(d(v)\cdot y)^2}\geq \frac{1}{2}\sum_{v'\in H^*\setminus C} \frac{(\ell_2(v,v'))^2}{(d(v)\cdot y)^2}
\end{align}
where $C$ is the set of vertices $v'$ for which $\ell_2^{B_i}(v,v')=1$.


\begin{clm}
$|H^*\setminus C|\leq|H^*|$
\end{clm}
\begin{proof}
The number of paths of length two originating at $v$ and going through $B_i$ is at least
$$\ell_2^{B_i}(v)\geq \eps^2 \cdot d(v)\cdot y/\log n\geq \frac{\eps^2 \cdot \theta\cdot \davg}{\log n}$$
(where $\ell_2^{B_i}(v)=\sum_{v'\in H^*}\ell_2^{B_i}(v,v')$). Also, $|H^*|\leq 2m/\theta$, so that for $\theta =2\sqrt{n\log n}/\eps$,
\begin{align}
|H^*|\leq \frac{\eps^2\sqrt{n}\davg}{2\sqrt{\log n}}  =\frac{\eps^2 \theta\davg}{2\log n}\leq \frac{1}{2}\ell_2^{B_i}(v).
\end{align}
Therefore,
\begin{align}
    \ell_2^{B_i, H^*\setminus C}(v)\geq \ell_2^{B_i}(v)-|H^*|\geq \frac{1}{2}\ell_2^{B_i}(v)
\end{align}

Hence,
\begin{align}
     \sum_{v'\in H^*\setminus C} \frac{(\ell_2^{B_i}(v,v'))^2}{(d(v)\cdot y)^2} \geq \frac{|H^*\setminus C|\cdot \left(\frac{\ell_2^{B_i, H*\setminus C}(v)}{|H^*\setminus C|}\right)^2}{(d(v)\cdot y)^2}
     \geq
     \frac{|H^*\setminus C|\cdot \left(\frac{\frac{1}{2}\ell_2^{B_i}(v)}{|H^*\setminus C|}\right)^2}{(d(v)\cdot y)^2}
     \geq
     \frac{(\eps d(v)\cdot y)^2}{|H^*\setminus C|\cdot (d(v)\cdot y)^2}\geq \frac{\eps^2}{|H^*|}
\end{align}

\end{proof}



\fi



\subsection{lower bound ideas for 4-cycle freeness}
Construction: a bipartite graph with sides $A\cup B$. $A$ is of size $\Theta(n)$ and all vertices have degree two. $B$ is of size $\Theta(\sqrt n)$ and with high probability all vertices have degree $\Theta^*(\sqrt n)$.
For the YES instances,  we have a perfect matching between vertices of $A$ and \emph{pairs of vertices} from $B$. We then add an edge between the matched vertex in $A$ and the two vertices in $B$. One can verify that this graph is 4-cycle free.
In the NO instances, we start as before, but then choose $O(1/\eps)$ of $A$'s vertices uar, and create copies of them and their incident edges. One can verify that each such copy creates a single 4-cycle, and hence we get a graph with $\Theta(\eps n)$ 4-cycles.

Intuition: in order to hit a vertex and it's copy one needs roughly $O(\sqrt n)$ samples by birthday paradox. Or from $H$'s side, one can hit a pair that belongs to a 4-cycle using $O(1/\eps)$ samples, but then these vertices have degree $\theta(\sqrt n)$ and it is not clear how to find their neighbors that belong to the 4-cycle.

\subsection*{Trying to deal with 4-cycles that have two  high-degree vertices (not both very high)}

As in the case where we had many $4$-cycles that contained two very-high degree vertices (i.e., in $H*$), we consider a bipartite graph where on one side are the low-degree vertices, $L$, and on the other side are the high-degree vertices, $H$ (and we may focus on such a graph after performing some manipulations that will cause use to pay $\poly(1/\eps)$ factors. We also assume that in $L$ we don't have vertices whose degree is significantly below the average degree (we need at some point to write something clean about these two things).

Suppose we buckets the vertices in $L$ as we did for the very-high degree case: $B_i=\{u\in L \mid d(u)\in [\eps\davg\cdot 2^{i-1}, \eps\davg\cdot 2^i \}$ (or a smaller factor than $2$ if needed). Then the collision probability is at least
\begin{equation}\sum_{v\in H} \frac{d(v)}{m} \sum_i \sum_{v'\in H} \frac{\ell_2^{B_i}(v,v')\cdot (\ell_2^{B_i}(v,v')-1)}{(d(v)\cdot 2^{i-1}\davg)^2} =
\sum_i \; \sum_{v\in H} \frac{d(v)}{m} \sum_{v'\in H} \frac{\ell_2^{B_i}(v,v')\cdot (\ell_2^{B_i}(v,v')-1)}{(d(v)\cdot 2^{i-1}\davg)^2}
\;.
\end{equation}
This implies that there exists a bucket $B_i$ such that the corresponding term in the sum over $i$ is at least $1/\log n$ the total collision probability.


\bibliographystyle{plainnat}
\bibliography{lit}

\end{document}
